% ============================================================
% Confusion Matrix — Hieroglyph Classification Example
% CSC491 Textbook, Chapter: Machine Learning
%
% A 2x2 grid showing TP, FP, FN, TN using the running
% hieroglyph example (classifying a glyph as ``Owl'').
%
% Preamble requirements:
%     \usepackage{tikz}
%     \usetikzlibrary{positioning}
%     \usepackage{xcolor}
%     \usepackage{fontspec}
%     \newfontfamily\egyptfont{Noto Sans Egyptian Hieroglyphs}
%
%     \definecolor{pipegreen}{RGB}{31, 120,  85}
%     \definecolor{barrierred}{RGB}{180, 40,  40}
%     \definecolor{pipenavy}{RGB}{26,  47,  89}
%     \definecolor{pipegold}{RGB}{160,110,   0}
% ============================================================

\begin{figure}[ht]
\centering
\begin{tikzpicture}[
  font=\sffamily\small,
  cell/.style={
    rectangle, rounded corners=3pt,
    minimum width=5.2cm, minimum height=3.0cm,
    align=center, text=white,
  },
]

% ── Axis labels ─────────────────────────────────────────────
\node[font=\sffamily\small\bfseries, text=pipenavy, rotate=90]
  at (-3.2, 0) {Actual Glyph};

\node[font=\sffamily\small\bfseries, text=pipenavy]
  at (2.8, 3.8) {Model's Prediction};

% Row labels (actual)
\node[font=\sffamily\small\bfseries, text=pipegreen, rotate=90, anchor=south]
  at (-2.6, 1.7) {Is an Owl};
\node[font=\sffamily\small\bfseries, text=barrierred, rotate=90, anchor=south]
  at (-2.6, -1.7) {Not an Owl};

% Column labels (predicted)
\node[font=\sffamily\small\bfseries, text=pipegreen]
  at (0.2, 3.35) {Predicted Owl};
\node[font=\sffamily\small\bfseries, text=barrierred]
  at (5.4, 3.35) {Predicted Not Owl};

% ── Quadrant cells ──────────────────────────────────────────

% True Positive (top-left) — correct identification
\node[cell, fill=pipegreen] (tp) at (0.2, 1.7)
  {\textbf{True Positive (TP)}\\[4pt]
   \footnotesize Model predicted {\egyptfont 𓅓} Owl\\
   \footnotesize \textit{Glyph was actually {\egyptfont 𓅓} Owl}\\[4pt]
   \scriptsize\color{white!75} Correct classification};

% False Negative (top-right) — missed owl
\node[cell, fill=pipegold!85!black] (fn) at (5.4, 1.7)
  {\textbf{False Negative (FN)}\\[4pt]
   \footnotesize Model predicted ``Not {\egyptfont 𓅓} Owl''\\
   \footnotesize \textit{Glyph was actually {\egyptfont 𓅓} Owl}\\[4pt]
   \scriptsize\color{white!75} Missed --- the {\egyptfont 𓅓} Owl was overlooked};

% False Positive (bottom-left) — misidentified as owl
\node[cell, fill=barrierred] (fp) at (0.2, -1.7)
  {\textbf{False Positive (FP)}\\[4pt]
   \footnotesize Model predicted {\egyptfont 𓅓} Owl\\
   \footnotesize \textit{Glyph was actually {\egyptfont 𓆑} Viper}\\[4pt]
   \scriptsize\color{white!75} Wrong --- {\egyptfont 𓆑} Viper mistaken for {\egyptfont 𓅓} Owl};

% True Negative (bottom-right) — correctly rejected
\node[cell, fill=pipegreen!70!black] (tn) at (5.4, -1.7)
  {\textbf{True Negative (TN)}\\[4pt]
   \footnotesize Model predicted ``Not {\egyptfont 𓅓} Owl''\\
   \footnotesize \textit{Glyph was indeed {\egyptfont 𓆑} Viper}\\[4pt]
   \scriptsize\color{white!75} Correct rejection};

\end{tikzpicture}
\caption{A confusion matrix for our hieroglyph classifier. Each quadrant shows what happens when the model's prediction is compared against the true label.}
\label{fig:confusion-matrix}
\end{figure}